\chaptersection{7}{Mitigation e Residual Risk}
A partire dalle minacce approfondite nel Capitolo 6, vengono individuate le principali misure di mitigazione applicabili ai relativi percorsi di attacco, collegando ciascuna contromisura alle condizioni e alle precondizioni emerse nei threat tree.

Le contromisure considerate nel seguito sono formulate come interventi di sicurezza applicabili ai percorsi di attacco individuati, con l'obiettivo di ridurre la probabilità di successo della minaccia, aumentarne la rilevabilità e, quando possibile, limitarne la propagazione lungo la pipeline. L'effetto di tali misure non implica tuttavia l'eliminazione completa del rischio: le condizioni che rimangono sfruttabili anche dopo l'applicazione delle contromisure vengono considerate nella successiva valutazione del rischio residuo.

\subsection{FM01.4 - Manipolazione intenzionale dell'immagine}
Il threat tree associato a FM01.4 individua quattro principali modalità di compromissione dell'immagine dermatologica: la fornitura diretta di un input avversariale, la modifica dell'immagine durante il trasferimento, la sua alterazione dopo l'acquisizione e prima dell'inferenza e l'inserimento di
contenuti visivi progettati per influenzare Qwen2-VL. 

Le mitigazioni devono pertanto agire sia sulla superficie di ingresso sia sulla protezione dell'immagine durante il suo percorso nella pipeline. Poiché un'immagine formalmente valida può comunque contenere perturbazioni o contenuti avversariali, i soli controlli sul formato del file non risultano sufficienti.
Per ciascun percorso vengono quindi considerate separatamente le principali precondizioni individuate nel threat tree, associando a ciascuna di esse una misura volta a impedirne lo sfruttamento o a renderlo più facilmente rilevabile.

{
\small
\setstretch{1}
\setlength{\tabcolsep}{4pt}
\renewcommand{\arraystretch}{1.07}

\begin{xltabular}{\textwidth}{
    >{\raggedright\arraybackslash}p{3.1cm}
    X
    X
}

\toprule

\textbf{Percorso di attacco} &
\textbf{Mitigazioni delle precondizioni} &
\textbf{Effetto atteso} \\

\midrule

\endfirsthead

\toprule

\textbf{Percorso di attacco} &
\textbf{Mitigazioni delle precondizioni} &
\textbf{Effetto atteso} \\

\midrule

\endhead

\multicolumn{3}{r}{\footnotesize Continua nella pagina successiva} \\
\endfoot

\endlastfoot


Fornire un'immagine avversariale all'ingresso &

1. Proteggere gli endpoint mediante API key o token autenticati e applicare
rate limiting.

2. Verificare la stabilità dell'inferenza rispetto a resize o ricompressione,
segnalando variazioni anomale di classe o confidence.

3. Decodificare e ricodificare lato server l'immagine nel formato e nelle
dimensioni attese prima dell'inferenza. &

Riduce la possibilità che un'immagine costruita intenzionalmente raggiunga
EfficientNet-B4 e Qwen2-VL senza essere rilevata e rende più difficile
mantenere perturbazioni efficaci dopo il preprocessing. \\

\midrule


Intercettare e modificare l'immagine durante il trasferimento &

1. Utilizzare TLS con verifica del certificato tra B4 e DermaTriage.

2. Associare all'immagine un HMAC-SHA-256 e verificarlo alla ricezione.

3. Vincolare il file ricevuto al relativo \texttt{consulto\_id} e verificarne
tipo e corretta decodifica. &

Riduce la possibilità di intercettare, sostituire o modificare l'immagine
durante il trasferimento e consente di rilevare alterazioni prima che il dato
venga utilizzato dalla pipeline. \\

\midrule


Modificare l'immagine dopo l'acquisizione e prima dell'inferenza &

1. Utilizzare un service account dedicato con permessi minimi sulle directory
temporanee.

2. Rendere l'immagine in sola lettura dopo l'acquisizione, limitando la
scrittura alla sola fase iniziale.

3. Calcolare un digest SHA-256 dopo l'acquisizione e verificarlo prima del
caricamento nei modelli. &

Riduce la possibilità che un soggetto con accesso all'ambiente di esecuzione
modifichi silenziosamente un'immagine già acquisita e permette di rilevare
alterazioni intervenute prima dell'inferenza. \\

\midrule


Inserire contenuti visivi avversariali destinati a Qwen2-VL &

1. Rilevare testo sovrapposto o elementi estranei al contenuto clinico atteso.

2. Vincolare Qwen2-VL tramite system prompt al solo compito di descrizione
clinica, trattando l'immagine come input non fidato.

3. Validare l'output e bloccare testo imperativo o campi estranei prima
dell'invio al RAG e a BioMistral. &

Riduce la probabilità che una visual prompt injection modifichi il
comportamento di Qwen2-VL e limita la propagazione di una descrizione alterata
verso il sistema RAG, BioMistral e gli stadi successivi. \\

\bottomrule

\caption{Mitigazioni proposte per FM01.4 in relazione alle precondizioni dei principali percorsi del threat tree}

\label{tab:mitigazioni-fm01-4}\\

\end{xltabular}
}

\subsection{FM13.1 - Contenuti avversariali nel feedback clinico}
Il threat tree associato a FM13.1 individua tre modalità attraverso cui un feedback clinico contaminato può influenzare l'aggiornamento della pipeline: l'inserimento di contenuti avversariali direttamente durante la revisione, l'alterazione di un feedback già registrato e la successiva inclusione del contenuto compromesso nel processo di evoluzione dei prompt.

Le mitigazioni devono quindi proteggere sia l'origine e la persistenza del feedback sia il suo utilizzo nel processo di aggiornamento. Particolare attenzione deve essere posta ai contenuti semanticamente avversariali ma clinicamente plausibili, che possono superare controlli esclusivamente sintattici e venire successivamente incorporati nelle nuove versioni dei prompt.

{
\small
\setstretch{1}
\setlength{\tabcolsep}{4pt}
\renewcommand{\arraystretch}{1.07}

\begin{xltabular}{\textwidth}{
    >{\raggedright\arraybackslash}p{3.1cm}
    X
    X
}

\toprule

\textbf{Percorso di attacco} &
\textbf{Mitigazioni delle precondizioni} &
\textbf{Effetto atteso} \\

\midrule

\endfirsthead

\toprule

\textbf{Percorso di attacco} &
\textbf{Mitigazioni delle precondizioni} &
\textbf{Effetto atteso} \\

\midrule

\endhead

\multicolumn{3}{r}{\footnotesize Continua nella pagina successiva} \\
\endfoot

\endlastfoot


Inserire un feedback avversariale all'origine &

1. Proteggere le funzioni di revisione mediante autenticazione nominativa,
MFA e autorizzazione RBAC per i ruoli abilitati.

2. Preferire campi strutturati e valori ammessi per le correzioni, limitando
il testo libero alle sole informazioni cliniche necessarie.

3. Sottoporre i feedback destinati all'aggiornamento a una seconda
validazione prima che possano essere utilizzati dal PromptManager. &

Riduce la possibilità che un soggetto non autorizzato o un singolo account
compromesso introduca contenuti avversariali e limita l'ingresso di feedback
manipolati nel processo di aggiornamento. \\

\midrule


Alterare un feedback già registrato &

1. Limitare l'accesso alle API e allo storage dei feedback mediante RBAC e
permessi di scrittura separati da quelli di sola consultazione.

2. Conservare le revisioni come versioni distinte, evitando la sovrascrittura
non tracciata del record clinico precedentemente validato.

3. Applicare HMAC ai record persistenti e mantenere audit log append-only
delle operazioni di modifica. &

Riduce la possibilità di modificare silenziosamente un feedback già validato
e consente di rilevare alterazioni prima del suo riutilizzo nei processi di
aggiornamento. \\

\midrule


Far includere il feedback contaminato nel processo di aggiornamento &

1. Includere nella finestra di aggiornamento soltanto feedback esplicitamente
approvati e non semplicemente le correzioni più recenti.

2. Separare il contenuto del feedback dalle istruzioni del processo di
evoluzione, trattandolo come dato non fidato all'interno del prompt.

3. Validare ogni nuova versione dei prompt su casi di riferimento prima
dell'attivazione, mantenendo versionamento e possibilità di rollback. &

Riduce la probabilità che un feedback contaminato venga incorporato nei prompt
attivi e limita la persistenza dell'alterazione nelle successive elaborazioni
di Qwen2-VL o BioMistral. \\

\bottomrule

\caption{Mitigazioni proposte per FM13.1 in relazione alle precondizioni dei principali percorsi del threat tree}

\label{tab:mitigazioni-fm13-1}\\

\end{xltabular}
}

\subsection{FM20.3 - Alterazione delle regole dell'Adaptation Layer}
Il threat tree associato a FM20.3 individua tre modalità attraverso cui un attaccante può compromettere la logica decisionale finale: la modifica diretta delle regole di mapping, l'alterazione dei parametri decisionali e il caricamento di una versione non autorizzata delle regole.

Le mitigazioni devono quindi proteggere sia il codice e la configurazione dell'Adaptation Layer sia il processo attraverso cui tali elementi vengono modificati o resi operativi. Poiché un'alterazione può produrre output formalmente validi ma clinicamente non coerenti con le regole previste, è necessario affiancare ai controlli di accesso anche verifiche di integrità, versionamento e validazione del comportamento decisionale.

{
\small
\setstretch{1}
\setlength{\tabcolsep}{4pt}
\renewcommand{\arraystretch}{1.07}

\begin{xltabular}{\textwidth}{
    >{\raggedright\arraybackslash}p{3.1cm}
    X
    X
}

\toprule

\textbf{Percorso di attacco} &
\textbf{Mitigazioni delle precondizioni} &
\textbf{Effetto atteso} \\

\midrule

\endfirsthead

\toprule

\textbf{Percorso di attacco} &
\textbf{Mitigazioni delle precondizioni} &
\textbf{Effetto atteso} \\

\midrule

\endhead

\multicolumn{3}{r}{\footnotesize Continua nella pagina successiva} \\
\endfoot

\endlastfoot


Modificare direttamente la logica di mapping &

1. Proteggere repository e componenti applicativi mediante RBAC, MFA e
permessi di modifica limitati ai soli ruoli autorizzati.

2. Richiedere code review e approvazione prima del merge di modifiche alla
funzione di mapping.

3. Eseguire test automatici che verifichino le corrispondenze
HIGH/MEDIUM/LOW-P1/P4 e la soglia di confidence prima del deploy. &

Riduce la possibilità di introdurre e rendere operativa una modifica non
autorizzata della logica e permette di rilevare variazioni rispetto alle regole
attese prima del rilascio. \\

\midrule


Alterare i parametri decisionali &

1. Separare i parametri decisionali dalla configurazione modificabile dagli
utenti e limitarne l'accesso in scrittura ai soli componenti autorizzati.

2. Vincolare la soglia di confidence e gli altri valori critici a configurazioni
versionate e sottoposte a controllo delle differenze.

3. Validare automaticamente la configurazione rispetto a valori ammessi e a
casi di test rappresentativi prima del caricamento. &

Riduce la possibilità di modificare silenziosamente i parametri utilizzati
dall'Adaptation Layer e rende rilevabili variazioni capaci di alterare
sistematicamente l'assegnazione della priorità. \\

\midrule


Caricare una versione non autorizzata delle regole &

1. Limitare l'accesso al meccanismo di deploy o aggiornamento mediante
autenticazione forte e privilegi amministrativi separati.

2. Firmare l'artefatto contenente le regole e verificarne la firma prima del
caricamento da parte del servizio.

3. Mantenere versionamento e hash della versione approvata, bloccando
l'attivazione di artefatti non corrispondenti a una release autorizzata. &

Riduce la possibilità di sostituire le regole con una versione alterata e
impedisce l'attivazione di artefatti privi di provenienza o integrità
verificabile. \\

\bottomrule

\caption{Mitigazioni proposte per FM20.3 in relazione alle precondizioni dei principali percorsi del threat tree}

\label{tab:mitigazioni-fm20-3}\\

\end{xltabular}
}

\subsection{FM21.5 - Esposizione non autorizzata dei dati clinici}

Il threat tree associato a FM21.5 individua tre modalità attraverso cui un
attaccante può ottenere accesso a informazioni cliniche riservate: la
compromissione dei meccanismi di autenticazione, lo sfruttamento di controlli
di autorizzazione insufficienti e l'accesso diretto alle risorse che
memorizzano o trasportano i dati.

Le mitigazioni devono quindi proteggere sia l'identità dei soggetti e dei
componenti che accedono ai servizi, sia l'autorizzazione sulle singole risorse
richieste, oltre ai dati conservati o trasferiti. Poiché un utente correttamente
autenticato può comunque tentare di accedere a consulti appartenenti ad altri
soggetti, autenticazione e autorizzazione devono essere trattate come controlli
distinti.

{
\small
\setstretch{1}
\setlength{\tabcolsep}{4pt}
\renewcommand{\arraystretch}{1.07}

\begin{xltabular}{\textwidth}{
    >{\raggedright\arraybackslash}p{3.1cm}
    X
    X
}

\toprule

\textbf{Percorso di attacco} &
\textbf{Mitigazioni delle precondizioni} &
\textbf{Effetto atteso} \\

\midrule

\endfirsthead

\toprule

\textbf{Percorso di attacco} &
\textbf{Mitigazioni delle precondizioni} &
\textbf{Effetto atteso} \\

\midrule

\endhead

\multicolumn{3}{r}{\footnotesize Continua nella pagina successiva} \\
\endfoot

\endlastfoot


Compromettere i meccanismi di autenticazione &

1. Proteggere gli account privilegiati mediante MFA e limitare l'uso di
credenziali condivise.

2. Applicare scadenza breve, rotazione e revoca di API key e JWT, evitando
token persistenti oltre il necessario.

3. Applicare rate limiting e rilevazione di tentativi anomali sugli endpoint
di autenticazione e sulle API che espongono dati clinici. &

Riduce la possibilità che credenziali o token compromessi vengano riutilizzati
con successo e limita l'accesso prolungato agli endpoint contenenti dati
clinici. \\

\midrule


Sfruttare controlli di autorizzazione insufficienti &

1. Verificare lato server, per ogni richiesta, la relazione tra identità
autenticata e \texttt{consulto\_id} richiesto.

2. Applicare autorizzazione object-level anche a immagini, documenti,
diagnosi e validazioni associate al consulto.

3. Rifiutare identificativi o parametri modificati che facciano riferimento
a risorse non appartenenti al soggetto autorizzato. &

Riduce la possibilità che un utente autenticato acceda a consulti o documenti
di altri soggetti mediante modifica degli identificativi o abuso dei propri
privilegi. \\

\midrule


Accedere direttamente ai dati o ai canali di trasferimento &

1. Cifrare database, backup e file persistenti contenenti dati clinici e
limitare l'accesso mediante account di servizio dedicati.

2. Proteggere file temporanei e log mediante permessi minimi, evitando la
registrazione di dati clinici non necessari.

3. Utilizzare TLS con verifica del certificato per le comunicazioni tra B4 e
DermaTriage e per gli altri trasferimenti contenenti dati sensibili. &

Riduce la possibilità di estrarre informazioni direttamente dallo storage,
dai file ausiliari o dal traffico di rete senza utilizzare i normali
meccanismi applicativi. \\

\bottomrule

\caption{Mitigazioni proposte per FM21.5 in relazione alle precondizioni dei principali percorsi del threat tree}

\label{tab:mitigazioni-fm21-5}\\

\end{xltabular}
}

\subsection{Valutazione del rischio residuo}
A seguito delle mitigazioni individuate, le quattro minacce vengono rivalutate utilizzando la stessa risk matrix qualitativa adottata nel Capitolo 5. Non viene quindi introdotto un nuovo metodo di scoring: Impact e Likelihood vengono riesaminati considerando l'effetto delle contromisure sui percorsi e sulle precondizioni individuate nei threat tree.

La rivalutazione considera sia la capacità delle misure preventive di rendere più difficile il completamento dell'attacco, sia i controlli di detection e recovery che possono limitarne gli effetti o impedirne la propagazione. Il livello residuo viene quindi ricavato combinando nuovamente Impact e Likelihood secondo la risk matrix già utilizzata nella prioritizzazione iniziale.

{
\small
\setstretch{1}
\setlength{\tabcolsep}{4pt}
\renewcommand{\arraystretch}{1.08}

\begin{xltabular}{0.75\textwidth}{
    >{\raggedright\arraybackslash}p{1.5cm}
    >{\raggedright\arraybackslash}p{2.4cm}
    >{\raggedright\arraybackslash}p{2.2cm}
    >{\raggedright\arraybackslash}p{2.5cm}
    >{\raggedright\arraybackslash}X
}

\toprule

\textbf{ID} &
\textbf{Rischio iniziale} &
\textbf{Impact residuo} &
\textbf{Likelihood residua} &
\textbf{Rischio residuo} \\

\midrule

\endfirsthead

\toprule

\textbf{ID} &
\textbf{Rischio iniziale} &
\textbf{Impact residuo} &
\textbf{Likelihood residua} &
\textbf{Rischio residuo} \\

\midrule

\endhead

\multicolumn{5}{r}{\footnotesize Continua nella pagina successiva} \\
\endfoot

\endlastfoot


FM01.4 &
Molto alto &
Alto &
Media &
Alto \\

\midrule

FM13.1 &
Alto &
Medio &
Bassa &
Basso \\

\midrule

FM20.3 &
Alto &
Alto &
Bassa &
Medio \\

\midrule

FM21.5 &
Alto &
Alto &
Bassa &
Medio \\

\bottomrule

\caption{Valutazione del rischio residuo delle minacce analizzate}

\label{tab:rischio-residuo}\\

\end{xltabular}
}

Per FM01.4 la Likelihood viene ridotta da Alta a Media. La protezione degli endpoint e dei canali di trasferimento, la normalizzazione dell'immagine e i
controlli applicati al ramo multimodale rendono più difficile completare i percorsi di attacco individuati. L'Impact rimane Alto, poiché un input
avversariale capace di superare tali controlli può ancora influenzare più stadi della pipeline e contribuire a una priorità clinica non corretta. Il rischio
residuo viene quindi valutato come Alto.

Per FM13.1 la Likelihood viene ridotta da Media a Bassa grazie ai controlli sull'inserimento, sulla modifica e sulla selezione dei feedback destinati
all'aggiornamento. L'Impact residuo viene valutato come Medio, poiché la validazione delle nuove versioni dei prompt, il versionamento e la possibilità
di rollback limitano la persistenza e la propagazione di un eventuale contenuto contaminato. Il rischio residuo viene pertanto valutato come Basso.

Per FM20.3 la Likelihood viene ridotta da Media a Bassa. La protezione del codice e dei parametri decisionali, insieme a code review, test automatici,
firma degli artefatti e controllo delle versioni, rende più difficile introdurre e attivare una logica di mapping alterata senza rilevazione.
L'Impact rimane Alto, poiché una compromissione riuscita può modificare sistematicamente la priorità clinica e le relative indicazioni operative. Il
rischio residuo viene quindi valutato come Medio.

Per FM21.5 la Likelihood viene ridotta da Media a Bassa. Autenticazione forte, autorizzazione sulle singole risorse, protezione dello storage e cifratura
delle comunicazioni riducono i principali percorsi di accesso non autorizzato ai dati clinici. L'Impact rimane Alto per la sensibilità delle informazioni e
per la possibile estensione dell'esposizione a più consulti. Il rischio residuo viene quindi valutato come Medio.

La valutazione evidenzia pertanto una riduzione del rischio per tutte le minacce considerate, senza assumere che le relative superfici di attacco
possano essere completamente eliminate. FM01.4 mantiene il livello residuo più elevato, mentre FM13.1 beneficia anche di misure di contenimento e ripristino
che ne riducono maggiormente gli effetti nel tempo.

\subsection{Assunzioni, limiti e controlli non pienamente verificabili}

L'analisi è sviluppata assumendo che l'architettura, i flussi di elaborazione,
le regole decisionali e i meccanismi di aggiornamento descritti nella
documentazione disponibile rappresentino la configurazione di riferimento di
DermaTriage. Sono inoltre considerate le interazioni documentate con B4 e i
componenti interni della pipeline, senza assumere la presenza di ulteriori
meccanismi autonomi, strumenti esterni o forme di memoria persistente non
descritte nell'architettura.

La documentazione consente di valutare i principali controlli e le superfici di
attacco a livello architetturale, ma non permette di verificare completamente
alcuni aspetti della configurazione operativa. Tra questi rientrano la
granularità effettiva delle policy di autorizzazione sulle singole risorse, le
modalità di rotazione e revoca delle credenziali, la configurazione concreta
dei canali TLS, la cifratura dello storage e dei backup, la protezione dei log,
l'isolamento dei processi e la configurazione del meccanismo di deploy.

Analogamente, non è possibile determinare dalla sola documentazione
l'effettiva applicazione di controlli quali MFA, firma degli artefatti,
verifiche HMAC, audit log append-only, separazione dei privilegi e procedure di
approvazione delle modifiche. Tali misure sono pertanto considerate nel piano
di mitigazione come controlli da applicare ai percorsi di attacco individuati,
mentre la loro efficacia effettiva richiederebbe una verifica della
configurazione distribuita e delle procedure operative adottate.

Restano inoltre fuori dal perimetro della presente analisi attività di verifica
dinamica quali penetration test, adversarial testing dei modelli, revisione
della configurazione runtime e test specifici sull'efficacia dei controlli di
monitoraggio e rilevazione. Tali attività potrebbero essere utilizzate per
validare sperimentalmente le ipotesi formulate nel threat model e affinare
ulteriormente la stima del rischio residuo.

Nel complesso, l'analisi identifica le principali superfici di attacco di
DermaTriage, ne valuta la priorità, ricostruisce i percorsi necessari alla loro
realizzazione e associa a ciascuno di essi contromisure coerenti con
l'architettura analizzata. La valutazione del rischio residuo completa infine
il processo distinguendo il rischio iniziale dalla criticità che permane dopo
l'applicazione delle mitigazioni considerate.